\begin{example}[Equal traces, different branching]\label{ex:branching}
Let the late-choice system start at $\ell_0$ and have exactly
\begin{equation}\label{eq:late-choice}
\ell_0\xrightarrow{a}\ell_1,\qquad
\ell_1\xrightarrow{b}\ell_b,\qquad \ell_1\xrightarrow{c}\ell_c,
\end{equation}
and let the early-choice system start at $e_0$ and have exactly
\begin{equation}\label{eq:early-choice}
e_0\xrightarrow{a}e_b,\quad e_0\xrightarrow{a}e_c,\quad
e_b\xrightarrow{b}e'_b,\quad e_c\xrightarrow{c}e'_c.
\end{equation}
All unlisted states are terminal; neither system has silent transitions or
cycles. Enumerating all paths, whose lengths are at most two, gives both exact
prefix-closed languages as $\{\epsilon,a,ab,ac\}$, not merely a truncated match.

An explicit simulation from early to late is
\begin{equation}\label{eq:early-late-witness}
R_{E\to L}=\{(e_0,\ell_0),(e_b,\ell_1),(e_c,\ell_1),
(e'_b,\ell_b),(e'_c,\ell_c)\}.
\end{equation}
Both initial $a$-moves are matched by $\ell_0\xrightarrow{a}\ell_1$.
At $(e_b,\ell_1)$ the sole $b$-move has a match; at $(e_c,\ell_1)$ the sole
$c$-move has a match. Terminal pairs impose no further obligations. Thus
$L_{\mathrm{early}}\weakpre L_{\mathrm{late}}$.

Conversely, suppose a simulation $R$ from late to early contains
$(\ell_0,e_0)$. Matching the single transition
$\ell_0\xrightarrow{a}\ell_1$ requires choosing one successor of $e_0$ and
including the resulting pair in $R$. If $(\ell_1,e_b)\in R$, the $c$-move
from $\ell_1$ has no match. If $(\ell_1,e_c)\in R$, its $b$-move has no
match. These exhaust the possible $a$-successors, so no such $R$ exists:
$L_{\mathrm{late}}\not\weakpre L_{\mathrm{early}}$. Weak bisimilarity
therefore also fails.

Different outgoing transitions from a related pair may indeed choose different
matching successors. However, the \emph{one} initial late $a$-transition must
be matched by an early successor pair that itself satisfies every subsequent
transfer obligation. Trace matching may select the $e_b$ branch for $ab$ and
the $e_c$ branch for $ac$ independently; simulation cannot combine those two
states into one defender state after the initial move. This is the quantifier
distinction behind the final strictness claim of Proposition
\ref{prop:hierarchy}.
\end{example}
